\section{Programme de formation sur 3 mois}
\label{sec:planning}

Le programme ci-dessous s'étale sur environ 3 mois, à raison d'une séance par semaine.
La première séance démarre volontairement par des \textbf{jeux de plage en PMT}
(palmes-masque-tuba, sans scaphandre) afin de mettre les élèves en confiance avant
l'introduction progressive du matériel de plongée.

\subsection{Vue d'ensemble}

\begin{longtable}{|c|p{2cm}|c|p{4.3cm}|p{5cm}|}
  \hline
  \rowcolor{ffessmblue}
  \textcolor{white}{\textbf{Séance}} & \textcolor{white}{\textbf{Semaine}} &
  \textcolor{white}{\textbf{Durée}} & \textcolor{white}{\textbf{Objectif principal}} &
  \textcolor{white}{\textbf{Compétences travaillées}} \\
  \hline
  \endfirsthead

  \hline
  \rowcolor{ffessmblue}
  \textcolor{white}{\textbf{Séance}} & \textcolor{white}{\textbf{Semaine}} &
  \textcolor{white}{\textbf{Durée}} & \textcolor{white}{\textbf{Objectif principal}} &
  \textcolor{white}{\textbf{Compétences travaillées}} \\
  \hline
  \endhead

  1 & S1 & 1h & Jeux de plage / accoutumance PMT &
    Se propulser ; Respecter le milieu ; Communiquer (bases) \\ \hline
  2 & S2 & 1h & S'équiper et se déséquiper &
    S'équiper et se déséquiper ; Se mettre à l'eau et sortir de l'eau (mise à l'eau simple) \\ \hline
  3 & S3 & 1h & Se mettre à l'eau et sortir de l'eau &
    Se mettre à l'eau et sortir de l'eau ; Se propulser (palmage ventral) ; Communiquer \\ \hline
  4 & S4 & 1h30 & S'immerger &
    Évoluer dans l'eau - S'immerger (canard, phoque) ; Se ventiler (tuba) \\ \hline
  5 & S5 & 1h30 & Ventilation en immersion et équilibrage &
    Se ventiler (immersion, vidage masque partiel) ; S'équilibrer (poumon ballast, initiation) \\ \hline
  6 & S6 & 1h30 & Vidage de masque et lâcher-reprise d'embout &
    Se ventiler (vidage masque complet, lâcher-reprise d'embout, apnée) \\ \hline
  7 & S7 & 1h30 & Gestion du gilet de stabilisation &
    S'équilibrer (gilet + poumon ballast combinés) \\ \hline
  8 & S8 & 1h30 & Palmage et communication avancés &
    Se propulser (toutes techniques) ; Communiquer (tous les signes) \\ \hline
  9 & S9 & 1h30 & Retourner en surface &
    Retourner en surface (vitesse, tour d'horizon, gonflage, remontée expiration contrôlée) \\ \hline
  10 & S10 & 1h30 & Évoluer en sécurité &
    Évoluer en sécurité (panne d'air, intervention en relais, essoufflement, crampe, malaise) \\ \hline
  11 & S11 & 1h30 & Mise en situation complète &
    Synthèse de toutes les compétences ; évaluation du palmage \\ \hline
  12 & S12 & 1h30 & Synthèse et évaluation intermédiaire &
    Bilan et reprise ciblée des compétences non acquises \\ \hline
  13 & S13 & 1h30 & Évaluation finale &
    Validation de l'ensemble des compétences du référentiel N1 \\ \hline
\end{longtable}

Les colonnes ``Semaine'' sont à faire correspondre aux dates réelles retenues par le club.

\subsection{Détail des séances}

\subsubsection*{Séance 1 --- Jeux de plage / accoutumance PMT (1h)}
\textbf{Compétences amorcées} : Se propulser, Respecter le milieu et l'environnement
(aisance aquatique), Communiquer (signes de base).

\begin{description}[leftmargin=2.2cm]
  \item[Échauffement (10 min)] Entrée libre dans l'eau, prise de contact avec le bassin,
    petits jeux collectifs en surface (ballon, chat perché).
  \item[Corps de séance (35 min)] Jeux de plage en palmes-masque-tuba : déplacements
    variés en palmage ventral et dorsal, passages sous une ligne d'eau, courtes apnées
    ludiques, récupération de petits objets au fond, relais palmés par équipes.
  \item[Retour au calme (15 min)] Étirements légers, échange sur les ressentis,
    présentation du programme de formation et découverte du matériel scaphandre.
  \item[Points de vigilance sécurité] Aucun signe d'essoufflement ou de stress ; adapter
    la profondeur du bassin ; s'assurer de l'aisance de chacun avant de complexifier les
    exercices.
\end{description}

\subsubsection*{Séance 2 --- S'équiper et se déséquiper (1h)}
\textbf{Compétences} : S'équiper et se déséquiper ; Se mettre à l'eau et sortir de l'eau
(mise à l'eau simple).

\begin{description}[leftmargin=2.2cm]
  \item[Échauffement (10 min)] Retour sur les sensations de la séance 1, mobilité
    articulaire.
  \item[Corps de séance (40 min)] Présentation du matériel (bouteille, gilet, détendeur),
    gréage et dégréage encadrés, vérification de la pression et du bon fonctionnement du
    gilet et du détendeur, capelage et décapelage au bord du bassin puis dans l'eau,
    première mise à l'eau simple en scaphandre, aisance en surface.
  \item[Retour au calme (10 min)] Dégréage, rinçage et rangement du matériel, rappel des
    règles d'hygiène et d'entretien.
  \item[Points de vigilance sécurité] Prévention des accidents liés aux chutes de
    bouteille et aux équipements sous pression ; toujours attacher le bloc ; vérifier le
    lestage adapté au milieu.
\end{description}

\subsubsection*{Séance 3 --- Se mettre à l'eau et sortir de l'eau (1h)}
\textbf{Compétences} : Se mettre à l'eau et sortir de l'eau ; Se propulser (palmage
ventral) ; Communiquer.

\begin{description}[leftmargin=2.2cm]
  \item[Échauffement (10 min)] Gréage encadré, avec autonomie croissante.
  \item[Corps de séance (40 min)] Saut droit, départ plage, sortie de l'eau (retrait de
    l'ensemble bloc-gilet en surface), palmage ventral en surface sur courte distance,
    introduction des signes OK / monter / descendre / ça ne va pas.
  \item[Retour au calme (10 min)] Dégréage, bilan collectif.
  \item[Points de vigilance sécurité] Hauteur de saut adaptée, vérifier l'absence
    d'autres plongeurs à l'entrée dans l'eau, communiquer avant toute mise à l'eau.
\end{description}

\subsubsection*{Séance 4 --- S'immerger (1h30)}
\textbf{Compétences} : Évoluer dans l'eau -- S'immerger (canard, phoque) ; Se ventiler
(ventilation et vidage du tuba).

\begin{description}[leftmargin=2.2cm]
  \item[Échauffement (10 min)] Rappel du palmage de surface, gréage.
  \item[Corps de séance (60 min)] Canard et phoque en PMT puis en scaphandre, recherche
    de l'équilibre à 3 m, ventilation en surface sur tuba et vidage du tuba.
  \item[Retour au calme (10 min)] Étirements, bilan.
  \item[Points de vigilance sécurité] Prévention des barotraumatismes de l'oreille :
    équilibrage dès la surface, lestage adapté au milieu et au matériel.
\end{description}

\subsubsection*{Séance 5 --- Ventilation en immersion et équilibrage (1h30)}
\textbf{Compétences} : Se ventiler (ventilation en immersion, vidage de masque
partiel) ; S'équilibrer (poumon ballast, initiation).

\begin{description}[leftmargin=2.2cm]
  \item[Échauffement (10 min)] Immersion et retour au calme respiratoire.
  \item[Corps de séance (65 min)] Maîtrise et régulation de la ventilation en immersion
    (fréquence, amplitude), introduction du vidage de masque partiel, initiation au
    poumon ballast pour varier légèrement la profondeur.
  \item[Retour au calme (15 min)] Bilan, échange sur les ressentis.
  \item[Points de vigilance sécurité] Progressivité stricte, aucune situation brutale,
    arrêt immédiat de l'exercice en cas de stress.
\end{description}

\subsubsection*{Séance 6 --- Vidage de masque et lâcher-reprise d'embout (1h30)}
\textbf{Compétences} : Se ventiler (vidage du masque complet, lâcher et reprise
d'embout, apnée).

\begin{description}[leftmargin=2.2cm]
  \item[Échauffement (10 min)] Rappel du vidage partiel de la séance 5.
  \item[Corps de séance (65 min)] Vidage de masque complet dans des situations variées,
    lâcher-reprise d'embout (vidage par expiration et utilisation du bouton de
    surpression), courtes apnées de déplacement (profondeur et distance modérées).
  \item[Retour au calme (15 min)] Bilan individuel.
  \item[Points de vigilance sécurité] Encadrant à proximité immédiate en permanence,
    respecter le rythme de chaque élève, jamais de situation générant de l'insécurité.
\end{description}

\subsubsection*{Séance 7 --- Gestion du gilet de stabilisation (1h30)}
\textbf{Compétences} : Évoluer dans l'eau -- S'équilibrer (gilet et poumon ballast
combinés).

\begin{description}[leftmargin=2.2cm]
  \item[Échauffement (10 min)] Rappel poumon ballast.
  \item[Corps de séance (65 min)] Utilisation de l'inflateur et des différentes purges,
    équilibrage en situation statique puis dynamique, combinaison gilet et poumon
    ballast, maintien d'une profondeur stable à $\pm$ 1 m.
  \item[Retour au calme (15 min)] Bilan.
  \item[Points de vigilance sécurité] Gonflage progressif du gilet, ne jamais provoquer
    de remontée incontrôlée par surgonflage.
\end{description}

\subsubsection*{Séance 8 --- Palmage et communication avancés (1h30)}
\textbf{Compétences} : Se propulser (toutes techniques de palmage) ; Communiquer
(ensemble des signes conventionnels).

\begin{description}[leftmargin=2.2cm]
  \item[Échauffement (10 min)] Palmage ventral et dorsal.
  \item[Corps de séance (65 min)] Palmage dorsal, palmage de sustentation, palmage en
    immersion, nage en capelé (simulation de retour au bateau), révision et complément
    des signes (mi-pression, réserve, panne d'air, essoufflement, froid, fin de
    plongée/d'exercice).
  \item[Retour au calme (15 min)] Bilan, gestion de l'effort.
  \item[Points de vigilance sécurité] Aucune recherche de performance chronométrée ;
    prévention de l'essoufflement.
\end{description}

\subsubsection*{Séance 9 --- Retourner en surface (1h30)}
\textbf{Compétences} : Retourner en surface (vitesse de remontée, tour d'horizon,
gonflage du gilet, remontée en expiration contrôlée, tenue d'un palier).

\begin{description}[leftmargin=2.2cm]
  \item[Échauffement (10 min)] Rappel de l'équilibrage.
  \item[Corps de séance (65 min)] Maîtrise de la vitesse de remontée avec et sans
    repères visuels, tour d'horizon et gonflage du gilet en surface, tenue d'un palier
    en pleine eau, remontée en expiration contrôlée depuis une profondeur n'excédant pas
    6 m.
  \item[Retour au calme (15 min)] Bilan.
  \item[Points de vigilance sécurité] Accent sur la prévention de la surpression
    pulmonaire : ne jamais bloquer sa ventilation pendant la remontée.
\end{description}

\subsubsection*{Séance 10 --- Évoluer en sécurité (1h30)}
\textbf{Compétences} : Évoluer en sécurité (procédures du GP, intervention en relais).

\begin{description}[leftmargin=2.2cm]
  \item[Échauffement (10 min)] Rappel des signes de communication.
  \item[Corps de séance (65 min)] Simulation de panne d'air (apnée expiratoire sur 10 m
    à l'horizontale puis utilisation de l'octopus du GP), intervention en relais auprès
    d'un équipier en difficulté (passage de l'octopus), réaction face à un essoufflement,
    une crampe, un malaise simulé.
  \item[Retour au calme (15 min)] Débriefing systématique après chaque simulation.
  \item[Points de vigilance sécurité] Exercices encadrés individuellement, procédures
    répétées calmement, jamais de simulation brutale.
\end{description}

\subsubsection*{Séance 11 --- Mise en situation complète (1h30)}
\textbf{Compétences} : synthèse de l'ensemble des compétences déjà travaillées ;
évaluation du palmage.

\begin{description}[leftmargin=2.2cm]
  \item[Échauffement (10 min)] Gréage autonome.
  \item[Corps de séance (65 min)] Parcours combinant équipement, mise à l'eau,
    immersion, déplacement, communication et retour en surface ; évaluation du palmage
    sur une distance de l'ordre de 100 m en PMT et 50 m en capelé (pas de chronométrage,
    seul critère : absence d'essoufflement).
  \item[Retour au calme (15 min)] Bilan collectif.
  \item[Points de vigilance sécurité] Veiller à la cohésion de la palanquée sur
    l'ensemble du parcours.
\end{description}

\subsubsection*{Séance 12 --- Synthèse et évaluation intermédiaire (1h30)}
\textbf{Objectif} : bilan individuel et reprise ciblée des compétences non acquises.

\begin{description}[leftmargin=2.2cm]
  \item[Corps de séance] Bilan individuel à partir de la fiche de suivi et d'évaluation
    (section~\ref{sec:fiche-evaluation}), reprise ciblée des compétences non acquises,
    exercices d'entraînement complémentaires selon les besoins de chaque élève.
  \item[Points de vigilance sécurité] Adapter le contenu au niveau réel de chaque élève,
    sans précipitation.
\end{description}

\subsubsection*{Séance 13 --- Évaluation finale (1h30)}
\textbf{Objectif} : validation de l'ensemble des compétences du référentiel N1.

\begin{description}[leftmargin=2.2cm]
  \item[Corps de séance] Évaluation finale de toutes les compétences pratiques,
    entretien oral sur les connaissances théoriques, et si possible première plongée en
    milieu naturel encadrée. Remise du bilan individuel à chaque élève.
  \item[Points de vigilance sécurité] Rappel : en cas de certification délivrée en
    milieu artificiel, au moins 4 plongées en milieu naturel devront être réalisées et
    attestées dans les 12 mois suivant l'obtention de la certification.
\end{description}
